\subsection{Exact wave isometry for the original residue forms}
\label{actual-residue-wave-isometry}

The Fourier estimate uses an isometry of the actual differential
forms, including their auxiliary components.  Multiplication by a
unit-modulus phase alone is insufficient for a differentiated form.
The prescribed component correction gives an exact isometry, with no
loss depending on the rank or the word length.

At one level of the original tower, let $\beta_w$ denote the
sesquilinear residue form for a finite signed word $w$.  We retain
the original compact, unipotent and terminal Haar measures, the gauge
$\prod_a|c_a|^{2(r-a)+k}$, and the outer measure
$\prod_a|c_a|^{-2}dc_a$.  The established local compatibility theorem
identifies this form with
\begin{equation}\label{wave-isometry-form}
 \beta_w(x,y)=C_r^{-1}\langle R_w\Phi_x,\Phi_y\rangle,
 \qquad C_r>0,
\end{equation}
on compact smooth coefficients with positive Gram cutoff and on the
corresponding smooth rapidly decreasing class.  The same $C_r$ works
for every finite word; it is not set to one
\cite{wave-isometry-established}.

For $s\in\mathbb C$ set
\[
 f_s(z)=\exp(sz-\bar s\bar z),\qquad
 U_s=\begin{pmatrix}1&0\\s&1\end{pmatrix}.
\]
At a plus letter the source component map is $U_s$, and at a
minus letter it is $\bar U_s$.  Let $V_s(w)$ be their tensor product
over the original word positions.  Its recursive action is precisely
the source's $\mathrm{uInv}(s,w)$: the component with new index zero
is unchanged, while the component with new index one gains $s$ times
the zero component at a plus letter and $\bar s$ times it at a
minus letter.  Define
\[
 (\mathcal W_{s,w}x)(g)
       =f_s(D_1(g))\,V_s(w)x(g),
\]
where $D_1$ is the first original leading minor at the current
level.  On the regular radial domain it is the first outer
coordinate, or the terminal coordinate $\zeta$ when the residual
unipotent rank is zero.

\begin{theorem}\label{wave-isometry-actual}
For every original admissible datum, every finite admissible signed
word $w$, every $s\in\mathbb C$, and all coefficient vectors $x,y$
in either class just specified,
\begin{equation}\label{wave-isometry-exact}
 \beta_w(\mathcal W_{s,w}x,\mathcal W_{s,w}y)=\beta_w(x,y).
\end{equation}
The wave preserves the same positive Gram cutoff and coefficient
class.  Thus it supplies both original interfaces
\texttt{VIso} and \texttt{VIsoCpt}, at all levels $j\le r$, with
their original complex phase and component convention
\cite{wave-isometry-source}.
\end{theorem}

\begin{proof}
We give the terminal calculation, the mixed-word covariance and
the analytic transfer separately.

At residual rank zero write $Z=\zeta$ and $B=\bar\zeta$ as
multiplication operators.  The recovered contracted matrix is
\begin{equation}\label{wave-isometry-contracted}
 Q(Y)=\sigma\begin{pmatrix}
 D&-[D,Z]\\ [D,B]&-[[D,Z],B]
 \end{pmatrix},\qquad \sigma\in\{1,-1\},
\end{equation}
where $D$ is the correctly shifted full Capelli determinant.
The spectral shifts and the real sign are unchanged.  Its standard
and dual coordinate laws are
$[E_{ab},\eta_i]=\delta_{bi}\eta_a$ and
$[E_{ab},y_j]=-\delta_{aj}y_b$, with
$Z=\sum_jx_j\eta_j$ and $B=\sum_i\xi_i y_i$ for the fixed
vectors $x,\xi$.  All coordinate multipliers commute.
The cofactor formula
$[D,Z]=\sum_{i,j}x_j\eta_i C_{ij}$, with the original shifted
cofactors, gives
\[
                         [[D,Z],Z]=0.
\]
Indeed the second commutator is a sum of complementary two-row,
two-column minors skew in the two omitted column indices, multiplied
by the symmetric coefficients $x_jx_l$.  This cancellation is valid
before commuting any minor past another operator.  The corresponding
row expansion gives $[[D,B],B]=0$, by skewness in the two indices
carrying $\xi$.  Jacobi's identity and $[Z,B]=0$ now show that
$[[D,Z],B]$ commutes with both $Z$ and $B$.

Let $M_s$ be multiplication by $f_s(Z)$.  Differentiating
$M_{-ts}DM_{ts}$ in the real variable $t$ and using the preceding
vanishing commutators gives the exact finite formula
\[
 M_{-s}DM_s=D+s[D,Z]-\bar s[D,B]
                              -|s|^2[[D,Z],B].
\]
The three other entries follow by taking commutators with $Z,B$.
Consequently
\begin{equation}\label{wave-isometry-matrix}
                 M_{-s}Q(Y)M_s=U_{-s}^*Q(Y)U_{-s}.
\end{equation}
This is a differential-operator identity on the common smooth
domain, so no exponential of an unbounded closed operator is involved.

At positive residual rank, the established one-pair restriction
expresses the initial matrix as $T(Y)^*Q_{\rm inner}(Y)T(Y)$,
with its original spectral shifts and
\[
 T(Y)=\begin{pmatrix}Y+p&-c\\c^{-1}&0\end{pmatrix},
 \qquad p=c\partial_c.
\]
The inner matrix is independent of $c$.  Since
$p f_s=f_s(p+sc)$,
\[
                  T(Y)M_s=M_sT(Y)U_{-s}.
\]
The formal adjoint for $|c|^{-2}dc$ has the corresponding identity.
Thus \eqref{wave-isometry-matrix} holds at this rank too.  All
remaining outer coordinates and their gauge factors are unchanged.

We next verify that \eqref{wave-isometry-matrix} survives every
mixed residue word.  This step cannot be replaced by treating a
partial transpose as an adjoint.  For an arbitrary invertible scalar
matrix $A$, replace an initial operator matrix $Q$ by $A^*QA$.
For each processed plus position put $A$ in that auxiliary factor,
and for each processed minus position put $\bar A$; denote the
resulting tensor product by $V$.  The current main matrix $Q_w(Y)$
and residue $R_w$ transform as
\begin{equation}\label{wave-isometry-word-covariance}
       Q_w(Y)\longmapsto V^{-1}A_0^*Q_w(Y)A_0V,
       \qquad R_w\longmapsto V^*R_wV.
\end{equation}
Here $0$ is the main index.  To prove this, process the source
list tail first, in its prescribed order.  At a plus step the
main matrix is multiplied on the right by
$((Y-y)I-\Pi_{0a})/(Y-y-1)$; at a minus step it is multiplied
on the left by $((y-Y)I-\Omega_{0a})/(y-Y-1)$.
The permutation and contraction identities
\[
 [\Pi_{0a},A_0A_a]=0,
 \qquad
 [\Omega_{0a},A_0^*\bar A_a^{-1}]=0
\]
prove the first transformation by induction.  The latter follows
from $\Omega=\sum e_{ij}\otimes e_{ij}$ and the invariance of
the evaluation tensor under $C\otimes C^{-t}$.
For the residue update, the old $V^{-1}$ cancels the old $V$
in $V^*R_wV$.  Inserting the new main matrix gives
$A_a^*(\operatorname{ins}_aQ_w)A_a$ at a plus step.
At a minus step, partial transposition of the main index gives
$A_a^t(\operatorname{ins}_aQ_w^t)\bar A_a$.
These are exactly the required new congruences.  No operator
coefficients are commuted, and no rational denominator changes.
This proves \eqref{wave-isometry-word-covariance} for every length.

Apply it with $A=U_{-s}$.  Multiplication by $M_s$ commutes with
all scalar auxiliary matrices, so the left transformation of the
residue is $M_{-s}R_wM_s$.  Also
$V_{-s}(w)V_s(w)=I$.  Since $|f_s|=1$, the original radial
inner product therefore satisfies
\[
 \begin{split}
 \langle R_wM_sV_s\Phi_x,M_sV_s\Phi_y\rangle
 &=\langle V_{-s}^*R_wV_{-s}V_s\Phi_x,V_s\Phi_y\rangle\\
 &=\langle R_w\Phi_x,\Phi_y\rangle.
 \end{split}
\]

Finally, $D_1(g)=h(gf'_0,f_0)$ is invariant under the original
integrating groups: each compact factor and each unipotent factor
fixes the corresponding first paired vector, which is the first
isotropic flag vector at positive rank and $z,z'$ at rank zero.
Thus $D_1(s^{-1}gt)=D_1(g)$ for the actual left and right
integrating elements $s,t$.  Hence the actual gauged Jacquet transform of the source
wave is $M_sV_s\Phi_x$.  Multiplication by the phase preserves
compact support and positive Gram support.  Every derivative of
the phase is the same phase times a polynomial in matrix entries;
for each fixed $s$ it therefore also preserves the class with
arbitrarily fast decay of all derivatives.  The constant auxiliary
shears preserve both classes.  All pairings and adjoint operations
above are consequently on the common classes of the established
compatibility theorem.  Multiplying the identity by the same
$C_r^{-1}$ on both sides proves \eqref{wave-isometry-exact}.
Its equality is uniform in $r,w,s$; no bound for a class-preservation
seminorm has been substituted for a form estimate.
\end{proof}

This closes the wave-isometry input of the actual kernel estimate.
The mixed determinant Cauchy inequality, the partial-head Fourier
identity on all relevant strata, and the final uniform normalizer
and rational-return comparisons remain separate obligations.
